\input{../tex-inputs/latex-korrekturansicht-vorspann}

\section[Arthur Schnitzler an Berta Zuckerkandl, 3. 6. 1931]{L03980 Arthur Schnitzler an Berta Zuckerkandl, 3. 6. 1931}
\nopagebreak\mylabel{L03980v}
\rehead{ }\normalsize\beginnumbering\briefempfaengerindex{Zuckerkandl, Berta@\textsc{Zuckerkandl, Berta}!zzzSchnitzler, Arthur@\emph{von Arthur Schnitzler}!1931-06-031@{3. 6. 1931}|(be}
\toendnotes[C]{\smallbreak\pagebreak[2]}
\correspDesc{Versand  durch Arthur Schnitzler am 3. 6. 1931 in Wien
\newline{}Erhalt  durch Berta Zuckerkandl im Zeitraum [4. 6. 1931
                  – 8. 6. 1931?] in Paris}\toendnotes[C]{\smallbreak}
\Standort{DLA, HS.1985.1.2282.}
\physDesc{Brief, Durchschlag, 1 Blatt, 1 Seite, 712 Zeichen
\newline{}Schreibmaschine
\newline{}Handschrift: roter Buntstift, lateinische Kurrent (\noindent{}beschriftet: »\uline{Zuckerkandl}« und »\uline{Paris}«, drei Unterstreichungen)}\toendnotes[C]{\smallbreak}
\pstart
           \raggedleft{}{\pb}3. 6. 1931.\pend
           
\pstart{}Liebe und verehrte Freundin.\pend\vspace{0.5em}
\pstart
           Dieser Tage kommt, wie Sie wissen, die »\textcolor{green}{Komödie der
                  Worte}\pwindex{Schnitzler, Arthur 15. 5. 1862 Wien – 21. 10. 1931 ebd.@\textsc{Schnitzler, Arthur} (15. 5. 1862 Wien – 21. 10. 1931 ebd.), \emph{Schriftsteller, Mediziner}!Komödie der Worte. Drei Einakter@\strich\emph{Komödie der Worte. Drei Einakter}|pw}{}\ledrightnote{\textcolor{green}{Komödie der Worte. Drei Einakter}}« in \textcolor{pink}{Paris}\oindex{Paris@\textbf{Paris}, \emph{Hauptstadt}|pw}{}\ledrightnote{\textcolor{pink}{Paris}}{ }\label{K_L03980-1v}\edtext{zur deutschen \textcolor{violet}{Aufführung}\eventindex{Théâtre des Champs-Élysées@\textbf{Théâtre des Champs-Élysées}!Premiere von Komödie der Worte auf Deutsch in Paris@Premiere von Komödie der Worte auf Deutsch in Paris|pwv}{}\ledrightnote{{$\rightarrow$}\emph{\textcolor{violet}{Premiere von Komödie der Worte auf Deutsch in Paris}}}}{\lemma{\textnormal{\emph{zur deutschen Aufführung}}}\Cendnote{\textnormal{Die \textcolor{pink}{schweizer}\oindex{Schweiz@\textbf{Schweiz}|pwk}
                  Regisseurin \textcolor{blue}{Georgette Boner}\pwindex{Boner, Georgette 4.\,2.\,1903 Mailand – 26.\,11.\,1998 Zürich@\textsc{Boner, Georgette} (4.\,2.\,1903 Mailand – 26.\,11.\,1998 Zürich), \emph{Regisseurin, Malerin, Literaturhistorikerin}|pwk}, die sich mit
                  einer Arbeit über \emph{\textcolor{green}{Arthur Schnitzlers
                     Frauengestalten}\pwindex{Boner, Georgette 4.\,2.\,1903 Mailand – 26.\,11.\,1998 Zürich@\textsc{Boner, Georgette} (4.\,2.\,1903 Mailand – 26.\,11.\,1998 Zürich), \emph{Regisseurin, Malerin, Literaturhistorikerin}!Arthur Schnitzlers Frauengestalten@\strich\emph{Arthur Schnitzlers Frauengestalten}|pwk}} promoviert hatte, inszenierte mit der Theatergruppe \emph{\textcolor{brown}{Studio Allemand}\orgindex{Studio Allemand@Studio Allemand|pwk}} den Einakterzyklus \emph{\textcolor{green}{Komödie der Worte}\pwindex{Schnitzler, Arthur 15. 5. 1862 Wien – 21. 10. 1931 ebd.@\textsc{Schnitzler, Arthur} (15. 5. 1862 Wien – 21. 10. 1931 ebd.), \emph{Schriftsteller, Mediziner}!Komödie der Worte. Drei Einakter@\strich\emph{Komödie der Worte. Drei Einakter}|pwk}}, der zwischen dem \textcolor{violet}{12. 6.}\eventindex{Théâtre des Champs-Élysées@\textbf{Théâtre des Champs-Élysées}!Premiere von Komödie der Worte auf Deutsch in Paris@Premiere von Komödie der Worte auf Deutsch in Paris|pwkv}
                  und dem 17. 06. 1931 am \textcolor{pink}{Théâtre des Champs-Élysées}\oindex{Théâtre des Champs-Élysées@\textbf{Théâtre des Champs-Élysées}|pwk} aufgeführt wurde.}}}\label{K_L03980-1}. Ich habe mir erlaubt
               sowohl Ihnen als Mme. \textcolor{blue}{Clemenceau}\pwindex{Clemenceau, Sophie 25.\,5.\,1862 – 24.\,9.\,1937@\textsc{Clemenceau, Sophie} (25.\,5.\,1862 – 24.\,9.\,1937)|pw}{}\ledrightnote{\textcolor{blue}{Sophie Clemenceau}} Billetts
               zusenden zu lassen und es wird mich freuen, wenn Sie Zeit haben sollten der \textcolor{violet}{Generalprobe}\eventindex{Paris@\textbf{Paris}!Generalprobe von Komödie der Worte auf Deutsch in Paris@Generalprobe von Komödie der Worte auf Deutsch in Paris|pwv}{}\ledrightnote{{$\rightarrow$}\emph{\textcolor{violet}{Generalprobe von Komödie der Worte auf Deutsch in Paris}}} beizuwohnen und
               mir vielleicht ein Wort darüber zu berichten. Indess hat Frau \textcolor{blue}{Clauser}\pwindex{Clauser, Suzanne 16.\,5.\,1898 Wien – 11.\,9.\,1981 Paris@\textsc{Clauser, Suzanne} (16.\,5.\,1898 Wien – 11.\,9.\,1981 Paris), \emph{Schriftstellerin, Übersetzerin}|pw}{}\ledrightnote{\textcolor{blue}{Suzanne Clauser}} auch eine ausgezeichnete \textcolor{green}{Uebersetzung}\pwindex{Schnitzler, Arthur 15. 5. 1862 Wien – 21. 10. 1931 ebd.@\textsc{Schnitzler, Arthur} (15. 5. 1862 Wien – 21. 10. 1931 ebd.), \emph{Schriftsteller, Mediziner}!?? [französische Übersetzung von Das Bacchusfest]@\strich\emph{?? [französische Übersetzung von Das Bacchusfest]}|pw}\pwindex{Schnitzler, Arthur 15. 5. 1862 Wien – 21. 10. 1931 ebd.@\textsc{Schnitzler, Arthur} (15. 5. 1862 Wien – 21. 10. 1931 ebd.), \emph{Schriftsteller, Mediziner}!?? [französische Übersetzung von Große Szene]@\strich\emph{?? [französische Übersetzung von Große Szene]}|pw}\pwindex{Schnitzler, Arthur 15. 5. 1862 Wien – 21. 10. 1931 ebd.@\textsc{Schnitzler, Arthur} (15. 5. 1862 Wien – 21. 10. 1931 ebd.), \emph{Schriftsteller, Mediziner}!?? [französische Übersetzung von Stunde des Erkennens]@\strich\emph{?? [französische Übersetzung von Stunde des Erkennens]}|pw}{}\ledrightnote{\textcolor{green}{?? [französische Übersetzung von Das Bacchusfest]}{\newline}\textcolor{green}{?? [französische Übersetzung von Große Szene]}{\newline}\textcolor{green}{?? [französische Übersetzung von Stunde des Erkennens]}} der drei \textcolor{green}{Einakter}\pwindex{Schnitzler, Arthur 15. 5. 1862 Wien – 21. 10. 1931 ebd.@\textsc{Schnitzler, Arthur} (15. 5. 1862 Wien – 21. 10. 1931 ebd.), \emph{Schriftsteller, Mediziner}!Bacchusfest@\strich\emph{Das Bacchusfest}|pwv}\pwindex{Schnitzler, Arthur 15. 5. 1862 Wien – 21. 10. 1931 ebd.@\textsc{Schnitzler, Arthur} (15. 5. 1862 Wien – 21. 10. 1931 ebd.), \emph{Schriftsteller, Mediziner}!Stunde des Erkennens@\strich\emph{Stunde des Erkennens}|pwv}\pwindex{Schnitzler, Arthur 15. 5. 1862 Wien – 21. 10. 1931 ebd.@\textsc{Schnitzler, Arthur} (15. 5. 1862 Wien – 21. 10. 1931 ebd.), \emph{Schriftsteller, Mediziner}!Große Szene@\strich\emph{Große Szene}|pwv}{}\ledrightnote{{$\rightarrow$}\emph{\textcolor{green}{Das Bacchusfest}}{\newline}{$\rightarrow$}\emph{\textcolor{green}{Stunde des Erkennens}}{\newline}{$\rightarrow$}\emph{\textcolor{green}{Große Szene}}} ins \textcolor{pink}{Französische}\oindex{Frankreich@\textbf{Frankreich}|pw}{}\ledrightnote{\textcolor{pink}{Frankreich}} fertig gestellt{[},{]}
               wir warten also damit, ebenso wie mit dem »\textcolor{green}{Bernhardi}\pwindex{Schnitzler, Arthur 15. 5. 1862 Wien – 21. 10. 1931 ebd.@\textsc{Schnitzler, Arthur} (15. 5. 1862 Wien – 21. 10. 1931 ebd.), \emph{Schriftsteller, Mediziner}!Professor Bernhardi. Komödie in fünf Akten@\strich\emph{Professor Bernhardi. Komödie in fünf Akten}|pw}\pwindex{Schnitzler, Arthur 15. 5. 1862 Wien – 21. 10. 1931 ebd.@\textsc{Schnitzler, Arthur} (15. 5. 1862 Wien – 21. 10. 1931 ebd.), \emph{Schriftsteller, Mediziner}!?? [französische Übersetzung von Professor Bernhardi]@\strich\emph{?? [französische Übersetzung von Professor Bernhardi]}|pw}{}\ledrightnote{\textcolor{green}{Professor Bernhardi. Komödie in fünf Akten}{\newline}\textcolor{green}{?? [französische Übersetzung von Professor Bernhardi]}}« und mit dem »\textcolor{green}{Weiten Land}\pwindex{Schnitzler, Arthur 15. 5. 1862 Wien – 21. 10. 1931 ebd.@\textsc{Schnitzler, Arthur} (15. 5. 1862 Wien – 21. 10. 1931 ebd.), \emph{Schriftsteller, Mediziner}!weite Land. Tragikomödie in fünf Akten@\strich\emph{Das weite Land. Tragikomödie in fünf Akten}|pw}{}\ledrightnote{\textcolor{green}{Das weite Land. Tragikomödie in fünf Akten}}«
               vor den Pforten der \textcolor{pink}{Pariser}\oindex{Paris@\textbf{Paris}, \emph{Hauptstadt}|pw}{}\ledrightnote{\textcolor{pink}{Paris}} Theater auf
               Einlass.\pend
           
\pstart
           Für heute nur dies und viele herzliche{\\[\baselineskip]} Grüsse. Hoffentlich habe ich bald
               die Freude Sie{\\[\baselineskip]} wiederzusehen.\pend
           \leftskip=0em{}{\vspace{1\baselineskip}}
\pstart
           \noindent{}Frau Hofrätin Bertha Zuckerkandl,{\\}\textcolor{pink}{Paris}\oindex{Paris@\textbf{Paris}, \emph{Hauptstadt}|pw}{}\ledrightnote{\textcolor{pink}{Paris}}.\pend
           \selectlanguage{ngerman}\endnumbering\briefempfaengerindex{Zuckerkandl, Berta@\textsc{Zuckerkandl, Berta}!zzzSchnitzler, Arthur@\emph{von Arthur Schnitzler}!1931-06-031@{3. 6. 1931}|)be}\mylabel{L03980h}
\begin{anhang}
\end{anhang}\newcommand{\dateiname}{L03980}\newcommand{\titel}{Arthur Schnitzler an Berta Zuckerkandl, 3. 6. 1931}\newcommand{\editorInnen}{Herausgegeben von Jahnke, SelmaMüller, Martin Anton}\input{../tex-inputs/latex-korrekturansicht-abspann}
